\section*{Bab \ref{ch:b1:titration-methods} --- Metode dan Kurva Titrasi}

\begin{solution}{exo:b1:titration-methods:1}
\ce{H3O+ + OH- -> 2H2O}, $K = 1/K_e = 10^{14.00}$. $C = 12.4 \times
0.100/20.0 = \qty{0.0620}{mol/L}$.
\end{solution}

\begin{solution}{exo:b1:titration-methods:2}
Pada \omterm{def:b1:titration-methods:titration}{ekuivalensi} larutannya berupa amonium klorida \qty{0.050}{mol/L}:
$\mathrm{pH} = \frac12(9.25 - \log 0.050) = 5.28$. Merah metil
(3.8--5.8) mencakupnya; jingga metil, yang trayeknya (2.5--4.5) baru
tercapai sesudah \omterm{def:b1:titration-methods:titration}{ekuivalensi} ketika pH turun, akan berubah terlambat.
\end{solution}

\begin{solution}{exo:b1:titration-methods:3}
\ce{MnO4- + 5Fe^2+ + 8H+ -> Mn^2+ + 5Fe^3+ + 4H2O};
$C = 5 \times 0.0200 \times 12.6/10.0 = \qty{0.126}{mol/L}$.
\end{solution}

\begin{solution}{exo:b1:titration-methods:4}
$C = 0.0100 \times 14.2/50.0 = \qty{2.84e-3}{mol/L}$ kalsium ditambah
magnesium. Keduanya bereaksi sebelum \omterm{def:b1:titration-methods:titration}{titik akhir}, yang hanya memberikan
jumlah keduanya: \omterm{def:b1:titration-methods:kinds}{titrasi simultan}.
\end{solution}

\begin{solution}{exo:b1:titration-methods:5}
0~mL: \omterm{def:b1:predominant-reaction:strong-acid}{asam lemah}, $\mathrm{pH} = \frac12(4.76 + 1.00) = 2.88$. 5.0~mL:
setengah \omterm{def:b1:titration-methods:titration}{ekuivalensi}, 4.76. 10.0~mL: 8.73
(\cref{prop:b1:titration-methods:ph-at-equivalence}). 15.0~mL: hidroksida
berlebih \qty{0.50}{mmol} dalam \qty{25.0}{mL}, $[\ce{OH-}] = 0.020$,
$\mathrm{pH} = 12.30$. Semuanya sesuai dengan kurva yang dihitung sampai
0.01.
\end{solution}

\begin{solution}{exo:b1:titration-methods:6}
Nilai $\mathrm{p}K_a$ 2.15 dan 7.21 berselisih lebih dari 4, demikian
pula 7.21 dan 12.34: dua \omterm{def:b1:titration-methods:titration}{ekuivalensi}, pada 10.0 dan \qty{20.0}{mL}.
Pertama: larutan \ce{H2PO4-}, $\mathrm{pH} \approx \frac12(2.15 + 7.21) = 4.68$;
kedua: \ce{HPO4^2-}, $\frac12(7.21 + 12.34) = 9.78$ (kurva eksaknya
memberikan 4.71 dan 9.66, karena rumus-rumus itu mengabaikan pengenceran
dan air). Keasaman ketiga ($\mathrm{p}K_a$ 12.34) terlalu dekat dengan
keasaman air: hidroksida yang ditambahkan hanya bereaksi sebagian dan pH
naik tanpa lonjakan.
\end{solution}

\begin{solution}{exo:b1:titration-methods:7}
Asam klorida: $6.2 \times 0.100/10.0 = \qty{0.062}{mol/L}$; asam etanoat:
$(14.6 - 6.2) \times 0.100/10.0 = \qty{0.084}{mol/L}$. Di tengah-tengah,
separuh asam etanoat telah dititrasi: $\mathrm{pH} = \mathrm{p}K_a = 4.76$.
\end{solution}

\begin{solution}{exo:b1:titration-methods:8}
$V_{\mathrm{eq}} = \qty{10.0}{mL}$. 2.0~mL: $E = 0.77 + 0.059\log(2/8) =
0.73$~V; 9.0~mL: $0.77 + 0.059\log 9 = 0.83$~V; 11.0~mL: permanganat
berlebih \qty{0.020}{mmol}, \ce{Mn^2+} \qty{0.200}{mmol}, $E = 1.51 +
\frac{0.059}{5}\log 0.10 = 1.50$~V; 20.0~mL: 1.51~V.
\end{solution}

\begin{solution}{exo:b1:titration-methods:9}
0~mL: $\kappa = (349.8 + 76.2) \times 0.0100 = \qty{4.26}{mS/cm}$.
10.0~mL: \ce{Na+} dan \ce{Cl-} pada $1.00/110 = \qty{9.09e-3}{mol/L}$,
$\kappa = (50.3 + 76.2) \times \num{9.09e-3} = 1.15$. 20.0~mL: \ce{Na+}
$0.0167$, \ce{Cl-} dan \ce{OH-} $0.00833$~mol/L, $\kappa = 0.84 + 0.635 + 1.65
= 3.13$~mS/cm. Kemiringannya (pengenceran diabaikan, \qty{1}{mL}
menambah \qty{1.0e-3}{mol/L}): $(50.3 - 349.8) \times 10^{-3} = -0.30$ dan
$(50.3 + 198.3) \times 10^{-3} = +0.25$~mS/cm per mL.
\end{solution}

\begin{solution}{exo:b1:titration-methods:10}
Pada $V_{\mathrm{eq}} - v$ asam berlebihnya $C_Bv$, pada $V_{\mathrm{eq}} + v$
basa berlebihnya $C_Bv$; dalam volume yang sama, $[\ce{H3O+}]$ sebelumnya
sama dengan $[\ce{OH-}]$ sesudahnya, sehingga $\mathrm{pH}(V_{\mathrm{eq}} - v) +
\mathrm{pH}(V_{\mathrm{eq}} + v) = 14$: simetris terhadap (V$_{\mathrm{eq}}$,
7). Lonjakan: \qty{1.0e-5}{mol} berlebih dalam \qty{20}{mL}, pH 3.30
sampai 10.70, 7.4 satuan. Jika dibagi 100: $\num{1.0e-7}$~mol dalam
\qty{20}{mL}, pH 5.3 sampai 8.7, 3.4 satuan: \omterm{def:b1:titration-methods:titration}{titik akhir} yang jauh kurang
tajam.
\end{solution}

\begin{solution}{exo:b1:titration-methods:11}
Sebelum: \ce{Fe^3+} yang terbentuk $= 5C_{\mathrm{Mn}}V$, \ce{Fe^2+} yang
tersisa $=
5C_{\mathrm{Mn}}(V_{\mathrm{eq}} - V)$, sehingga $E = 0.77 + 0.059\log\frac{V}{V_{\mathrm{eq}}
- V}$: pada $V_{\mathrm{eq}}/2$, $E = 0.77$. Sesudah: \ce{MnO4-} berlebih $\propto
V - V_{\mathrm{eq}}$, \ce{Mn^2+} $\propto V_{\mathrm{eq}}$: $E = 1.51 +
\frac{0.059}{5}\log\frac{V - V_{\mathrm{eq}}}{V_{\mathrm{eq}}}$ (pH 0); pada
$2V_{\mathrm{eq}}$, 1.51~V. Pada \omterm{def:b1:titration-methods:titration}{ekuivalensi} $(0.77 + 5 \times 1.51)/6 =
1.39$~V. Pada pH 1 pasangan permanganat lebih rendah sebesar $0.059 \times 8/5 =
0.094$~V: $E_{\mathrm{eq}} = (0.77 + 5 \times 1.416)/6 = 1.31$~V.
\end{solution}

\begin{solution}{exo:b1:titration-methods:12}
\ce{NH4+ + OH- -> NH3 + H2O} mempunyai $K = 10^{14.00 - 9.25} = 10^{4.75}$:
reaksinya tidak sempurna di dekat \omterm{def:b1:titration-methods:titration}{ekuivalensi} dan lonjakannya kecil, di
sekitar pH 11 (pH \omterm{def:b1:titration-methods:titration}{ekuivalensi} $14 - \frac12(4.75 + 1.30) = 11.0$). \omterm{def:b1:titration-methods:kinds}{Titrasi
balik}: tambahkan $n_1$ natrium hidroksida (berlebih), didihkan untuk
mengusir amonia yang terbentuk, lalu \omterm{def:b1:titration-methods:titration}{titrasi} hidroksida yang tersisa
dengan asam klorida ($n_2$, lonjakan kuat--kuat yang tajam):
$n(\ce{NH4+}) = n_1 - n_2$.
\end{solution}

\begin{solution}{pb:b1:titration-methods:1}
\textbf{1.} \ce{C6H6O6 + 2H+ + 2e- -> C6H8O6}.
\textbf{2.} \ce{C6H8O6 + I2 -> C6H6O6 + 2I- + 2H+}.
\textbf{3.} 0 dan $-$I.
\textbf{4.} \ce{I2 + 2S2O3^2- -> 2I- + S4O6^2-}, $\log K = 2(0.53 - 0.02)/0.059
= 17$.
\textbf{5.} Warna biru memerlukan iodin; warna itu hilang bersama iodin
terakhir, pada \omterm{def:b1:titration-methods:titration}{ekuivalensi}.
\textbf{6.} Kehilangan karena udara akan menurunkan hasilnya; jika
ditambahkan sekaligus dan berlebih, iodin mengoksidasi asam askorbat
dengan cepat dan sempurna.
\textbf{7.} Reaksi yang cepat dan \omterm{def:b1:extent-q-and-k:quantitative}{kuantitatif} selama seluruh \omterm{def:b1:titration-methods:titration}{titrasi},
tanpa oksidasi oleh udara sementara \omterm{def:b1:titration-methods:titration}{titran} ditambahkan perlahan-lahan,
dan \omterm{def:b1:titration-methods:titration}{titik akhir} pada kelebihan iodin yang pertama.
\textbf{8.} Iodin berlebih $n_R$; reaksi dengan asam askorbat; \omterm{def:b1:titration-methods:titration}{titrasi}
iodin yang tersisa dengan tiosulfat.
\textbf{9.} Jika tidak, sebagian asam askorbat akan tersisa dan kelebihan
iodin, yang nol, tidak akan menunjukkan berapa banyak. Di sini
\qty{2.775e-4}{mol} bereaksi dari \qty{5.000e-4}{mol}: memang berlebih.
\textbf{10.} Larutan tiosulfat perlahan-lahan berubah seiring waktu;
konsentrasinya masuk langsung ke dalam hasil.
\textbf{11.} 10.
\textbf{12.} Pipet volumetrik \qty{20.00}{mL}.
\textbf{13.} $20.00 \times 10^{-3} \times 0.0250 = \qty{5.000e-4}{mol}$.
\textbf{14.} $8.90 \times 10^{-3} \times 0.0500 = \qty{4.450e-4}{mol}$.
\textbf{15.} Separuhnya: \qty{2.225e-4}{mol}.
\textbf{16.} $5.000 - 2.225 = \qty{2.775e-4}{mol}$.
\textbf{17.} Satu banding satu: \qty{2.775e-4}{mol}.
\textbf{18.} $\times 10$: \qty{2.775e-3}{mol}, $\times 176.0$:
\qty{0.488}{g}.
\textbf{19.} Sekitar \qty{2}{\percent} di bawah label.
\textbf{20.} $n = \bigl(C_IV_I - \frac12C_TV_T\bigr)\,V_{\mathrm{labu}}/V_{\mathrm{alikuot}}$.
\textbf{21.} Iodin yang dimasukkan: pipet $0.03/20.00 = \qty{0.15}{\percent}$
dan konsentrasi \qty{0.2}{\percent}, sekitar \qty{0.25}{\percent}, yaitu
\qty{1.3e-6}{mol}. Kelebihan: buret $0.05/8.90 = \qty{0.56}{\percent}$ dan
konsentrasi \qty{0.2}{\percent}, sekitar \qty{0.60}{\percent}, yaitu
\qty{1.3e-6}{mol}.
\textbf{22.} Ketidakpastian mutlaknya bergabung, sedangkan selisihnya
lebih kecil daripada masing-masing jumlah: $\sqrt{1.3^2 + 1.3^2} \times 10^{-6} =
\qty{1.8e-6}{mol}$, \qty{0.66}{\percent} dari \qty{2.775e-4}{mol}, lebih
besar daripada masing-masing suku.
\textbf{23.} Dengan labu (\qty{0.1}{\percent}) dan pipet alikuot
(\qty{0.2}{\percent}): $\sqrt{0.66^2 + 0.1^2 + 0.2^2} = \qty{0.70}{\percent}$.
\textbf{24.} $0.488 \pm 0.003$~g: \textbf{\qty{0.49}{g} asam askorbat
per tablet}.
\end{solution}
